\section{Spectral Certification of the Reference Phase and G6}
\label{app:g6-certification}

This appendix supplies the algebra behind
Proposition~\ref{prop:reference-edge},
Lemma~\ref{lem:c6-realization}, and
Proposition~\ref{prop:g6-positive-completeness}. We first construct a
bilateral eigenfunction through the unsquared stable--unstable matching
condition. Elimination then identifies its squared energy, and the original
matching equation excludes the additional algebraic candidates.

Throughout, all interval endpoints are rational.  An assertion that a
polynomial or a radical expression has a fixed sign on an interval means
that repeated squaring, with the sign of each unsquared side retained,
gives the displayed rational enclosure.  Sturm counts use the ordinary
signed remainder sequence over \(\mathbb Q\).

\subsection{The reference Floquet calculation}
\label{appsubsec:reference-floquet}

The reference lift is
\[
 \tau_{\mathrm{ref}}=(1,1,-1,1,-1,-1,1,-1).
\]
For \(i=8m+r\), put \(u_i=z^m v_r\).  In the ordered basis
\((v_0,\ldots,v_7)\), the unsquared Bloch fiber is
{\small
\[
A_{\mathrm{ref}}(z)=
\begin{pmatrix}
0&1&1&0&0&0&z^{-1}&z^{-1}\\
1&0&1&1&0&0&0&-z^{-1}\\
1&1&0&1&-1&0&0&0\\
0&1&1&0&1&1&0&0\\
0&0&-1&1&0&1&-1&0\\
0&0&0&1&1&0&1&-1\\
z&0&0&0&-1&1&0&1\\
z&-z&0&0&0&-1&1&0
\end{pmatrix}.
\]
}
Reversing every undirected transition gives
\(A_{\mathrm{ref}}(z)^{\mathsf T}=A_{\mathrm{ref}}(z^{-1})\), so the
fiber is Hermitian for \(|z|=1\).  Fraction-free expansion gives
\begin{align}
\det(xI-A_{\mathrm{ref}}(z))={}&x^8-16x^6
 +(80-2(z+z^{-1}))x^4 \notag\\
&+(-128+16(z+z^{-1}))x^2
 +(z^2+z^{-2})-13(z+z^{-1})+40.       \label{eq:app-ref-det}
\end{align}
With \(y=x^2\) and \(c=z+z^{-1}\), this becomes
\begin{equation}
 P(y,c)=y^4-16y^3+(80-2c)y^2+(-128+16c)y
        +c^2-13c+38.                 \label{eq:app-ref-P}
\end{equation}
On the unit circle, \(c\in[-2,2]\).  At \(c=2\), translation by four in
the spectral variable gives
\[
 P(X+4,2)=X^4-20X^2+80.
\]
Consequently the upper squared edge at \(z=1\) is
\[
 \eta=4+s,
 \qquad s=\sqrt{10+2\sqrt5}.
\]
To prove that no other fiber reaches it, put \(u=y-\eta\) and \(t=2-c\).
Direct substitution in~\eqref{eq:app-ref-P} yields
\begin{align}
P(\eta+u,2-t)={}&u^4+4su^3+2u^2t
 +(40+12\sqrt5)u^2+4sut \notag\\
&+8\sqrt5\,su+t^2+(4\sqrt5-3)t.      \label{eq:app-edge-positive}
\end{align}
Every coefficient is positive.  For \(u,t\geq0\), the right-hand side
vanishes only at \((u,t)=(0,0)\).  Hence \(y\leq\eta\) on every unit-circle
fiber, and equality forces \(c=2\), hence \(z=1\).  This proves both the
value and the uniqueness assertion in
Proposition~\ref{prop:reference-edge}.  It also gives the useful strict
bound \(\eta<391/50<8\).

The same algebra appears in transfer form.  For a general lift \(\tau\),
write
\[
 X_i=(u_{i+1},u_i,u_{i-1},u_{i-2})^{\mathsf T}.
\]
The equation \(A_\tau u=\lambda u\) is equivalent to
\begin{equation}
 X_{i+1}=T_i(\lambda)X_i,
 \qquad
 T_i(\lambda)=
 \begin{pmatrix}
 -\tau_i&\tau_i\lambda&-\tau_i&-\tau_i\tau_{i-2}\\
 1&0&0&0\\0&1&0&0\\0&0&1&0
 \end{pmatrix}.                       \label{eq:app-step-transfer}
\end{equation}
Each \(T_i\) is invertible.  For one reference cell set
\(M_8(\lambda)=T_7(\lambda)\cdots T_0(\lambda)\). Multiplication
gives, with \(y=\lambda^2\),
\begin{equation}
 \det(zI-M_8(\lambda))
 =z^4+a(y)z^3+b(y)z^2+a(y)z+1,         \label{eq:app-palindromic}
\end{equation}
where
\begin{equation}
 a(y)=-2y^2+16y-13,\qquad
 b(y)=y^4-16y^3+80y^2-128y+40.         \label{eq:app-ab}
\end{equation}
After division by \(z^2\), the substitution \(w=z+z^{-1}\) reduces the
quartic to
\begin{equation}
 f(w;y)=w^2+a(y)w+b(y)-2=0.             \label{eq:app-w-quadratic}
\end{equation}
Its discriminant is \(-12y^2+96y+17\), and the reciprocal multiplier pairs
are obtained from
\begin{equation}
 w_\mp(y)=\frac{2y^2-16y+13
 \mp\sqrt{-12y^2+96y+17}}2,\qquad
 z_{s,u}(w)=\frac{w\mp\sqrt{w^2-4}}2.  \label{eq:app-multipliers}
\end{equation}

Equation~\eqref{eq:app-edge-positive} excludes a unit-circle multiplier for
every \(y>\eta\).  Reciprocity and continuity then give two multipliers in
\(|z|<1\) and two in \(|z|>1\), counted algebraically. The
multiplier discriminant factors as
\begin{align}
 &(y^2-12y+34)(y^2-4y+2)(12y^2-96y-17)^2 \notag\\
 &\hspace{35mm}{}\cdot
 (y^4-16y^3+76y^2-96y+16).              \label{eq:app-mult-discriminant}
\end{align}
Above the upper endpoint of \(J_6\), defined below, the only collision is
\begin{equation}
 y_*=4+\frac{\sqrt{627}}6,\qquad w=\frac{95}{12}>2.  \label{eq:app-repeat}
\end{equation}
Indeed, the positive roots of the first two quadratic factors in
\eqref{eq:app-mult-discriminant} satisfy
\[
6+\sqrt2<\frac{15}{2}<a_6,\qquad
2+\sqrt2<\frac72<a_6,
\]
and every root of the final quartic is at most \(\eta<a_6\) by the
reference-edge calculation.  The factor \(12y^2-96y-17\) has one negative
root and the single positive root \(y_*\). Thus \(y_*\) is the unique
collision on
\([a_6,16]\).
It is therefore a collision inside the hyperbolic region, not a loss of the
stable--unstable splitting.

\subsection{Transfer through the G6 core}
\label{appsubsec:g6-transfer}

For the forward interface, take
\[
 Q_i=1
 \quad\Longleftrightarrow\quad
 \bigl(i\leq0,\ i\equiv0\pmod4\bigr)
 \ \hbox{or}\
 \bigl(i\geq6,\ i\equiv6\pmod4\bigr),
\]
fix \(\tau_0=1\), and impose \(\tau_{i+1}=Q_i\tau_i\).  The left and
right cuts are \(-8\) and \(14\).  Thus
\begin{equation}
 D_6(\lambda)=T_{13}(\lambda)\cdots T_{-8}(\lambda),
 \qquad \det D_6(\lambda)=1.            \label{eq:app-D6}
\end{equation}
Equations~\eqref{eq:app-step-transfer} and~\eqref{eq:app-D6}, together with
the displayed definition of \(Q\), are a finite specification of
every entry of the defect transfer; each entry is an integer polynomial in
\(\lambda\) of degree at most eleven.  The two bulk monodromies at these
cuts are
\[
 M_L=T_{-9}\cdots T_{-16},\qquad
 M_R=T_{21}\cdots T_{14}.
\]
They are transfer representatives of the translated reference phases
\(B_0\) and \(B_2\).

For \(y>\eta\) and \(\lambda=\sqrt y>0\), let \(U_L(\lambda)\) be the
two-dimensional space of left-cut states decaying under backward iteration,
and let \(S_R(\lambda)\) be the two-dimensional space of right-cut states
decaying under forward iteration.  Since the recurrence is invertible, a
nonzero bilateral \(\ell^2\) solution exists if and only if
\begin{equation}
 D_6(\lambda)U_L(\lambda)\cap S_R(\lambda)\neq\{0\}.
                                                        \label{eq:app-physical-match}
\end{equation}
If \(u_1,u_2\) and \(s_1,s_2\) are bases of these planes, define
\begin{equation}
 \mathcal E_6(y)=
 \det[D_6(\sqrt y)u_1,D_6(\sqrt y)u_2,s_1,s_2].       \label{eq:app-Evans}
\end{equation}
A change of basis multiplies \(\mathcal E_6\) by a nowhere-zero factor.
Thus its zero set is precisely the coordinate-free condition
\eqref{eq:app-physical-match}.  In particular, changing cuts by full
periods or changing Grassmann charts cannot create or remove a physical
zero.

For the signed interval calculations below we fix the normalization
completely.  Order the two stable multipliers by
\[
0<z_1<z_2<1
\]
and the unstable multipliers as \(z_1^{-1},z_2^{-1}\).  For a row set
\(a<b<c\), define the cofactor vector by
\[
\bigl(v_{abc}(z)\bigr)_j
=(-1)^j\det\bigl((M-zI)_{\{a,b,c\},\,\{0,1,2,3\}\setminus\{j\}}\bigr).
\]
Use rows \(012\), divide each two-vector wedge by the ordered Vandermonde
\(z_1-z_2\), and order the four columns in
\eqref{eq:app-Evans} as
\[
\bigl[D_6u(z_1^{-1}),D_6u(z_2^{-1}),s(z_1),s(z_2)\bigr].
\]
This convention fixes a continuous orientation on \(J_6\).  The endpoint
and derivative signs below refer to this normalized determinant, rather than
to an arbitrary basis representative of the same zero set.

The essential-spectrum assertion used here can also be seen directly.
Cutting \(H_6=A_6^2\) on both sides of its finite core changes it by finite
rank into two periodic half-line compressions and a finite block.  A cutoff
Bloch wave proves that the spectrum of a whole-line periodic operator lies
in the essential spectrum of its half-line compression.  Conversely, if
\(R=(H_{\mathrm{per}}-x)^{-1}\), then the compressed operator \(PRP\) is a
two-sided inverse modulo finite-rank cross-boundary terms.  Hence the
half-line compression is Fredholm whenever \(x\) is outside the whole-line
spectrum.  It follows that
\[
 \sigma_{\mathrm{ess}}(H_6)=\sigma(H_{\mathrm{ref}}),
 \qquad \sup\sigma_{\mathrm{ess}}(H_6)=\eta.
\]
For \(y>\eta\), the hyperbolic splitting above shows that every matched
state decays exponentially.  This proves that
\eqref{eq:app-physical-match} is both necessary and sufficient for a
discrete physical level.

\subsection{Physical realization and algebraic identification}
\label{appsubsec:c6-identification}

Put
\[
 a_6=\frac{7905369311620327}{10^{15}},\qquad
 b_6=\frac{7905369311620328}{10^{15}}
     =\frac{988171163952541}{125000000000000},
 \qquad J_6=[a_6,b_6].
\]
On \(J_6\), the four multipliers are distinct and positive, with two on
each side of the unit circle.  Cofactor eigenvectors formed from rows
\(0,1,2\) of \(M-zI\) have a nonzero component throughout this interval,
so they define a valid normalization of~\eqref{eq:app-Evans}. Certified
rational interval evaluation gives
\[
 \mathcal E_6(a_6)<0<\mathcal E_6(b_6),
 \qquad
 \inf_{y\in J_6}\mathcal E_6'(y)>0.
\]
The explicit outward enclosures are recorded in the supplement. The
derivative is with respect to \(y\). The intermediate-value theorem and
strict monotonicity therefore prove that there is exactly
one positive-branch physical zero in \(J_6\).  Since \(y>0\), it is also a
simple zero as a function of \(\lambda\).  If the matching intersection had
dimension at least two, the matching matrix in~\eqref{eq:app-Evans} would
have rank at most two; its adjugate and hence the first derivative of its
determinant would vanish.  Therefore the positive eigenspace is
one-dimensional.

It remains to identify the squared energy.  Let \(z_1,z_2\) be the stable
multipliers and put
\[
 S=z_1+z_2,\qquad P=z_1z_2.
\]
The reciprocal quartic~\eqref{eq:app-palindromic} gives
\begin{equation}
 S(P+1)+a(y)P=0,\qquad
 P^2+S^2+1-b(y)P=0.                  \label{eq:app-SP-relations}
\end{equation}
After dividing the stable wedge by the nonzero Vandermonde
\(z_1-z_2\), every Pl\"ucker coordinate is symmetric in \(z_1,z_2\) and
hence polynomial in \(S,P\).  Substituting the first relation in the
second gives the product equation
\begin{equation}
 R(P,y):=(P^2+1-b(y)P)(P+1)^2+a(y)^2P^2=0.            \label{eq:app-product-relation}
\end{equation}
After removal of the proved nonzero chart denominators, the unsquared
matching numerator has the form
\[
 E_0(y,P)+\lambda E_1(y,P).
\]
This is also a finite specification of \(E_0,E_1\): form the matrices
\(T_i\) in~\eqref{eq:app-step-transfer}, multiply them in the order
\eqref{eq:app-D6}, form the signed cofactor vectors just defined, divide the
wedges by their ordered Vandermonde factors, and reduce the resulting
symmetric polynomials by~\eqref{eq:app-SP-relations}.
Thus every physical zero annihilates
\begin{equation}
 G(P,y)=E_0(y,P)^2-yE_1(y,P)^2.                       \label{eq:app-squared-Evans}
\end{equation}
The resultant of \(R\) and \(G\) with respect to \(P\) has degree 108 in
\(y\). Up to a nonzero rational constant, its factorization is
\begin{equation}
 r_2(y)^{16}r_{4,1}(y)^2r_{4,2}(y)^2p_6(y)^2q_6(y)^4, \label{eq:app-resultant-factorization}
\end{equation}
where
\begin{align*}
r_2(y)={}&2y^2-16y+13,\\
r_{4,1}(y)={}&y^4-16y^3+68y^2-32y+52,\\
r_{4,2}(y)={}&9y^4-144y^3+708y^2-1056y+404,\\
p_6(y)={}&16y^{10}-520y^9+6913y^8-48448y^7+191768y^6\\
&-423904y^5+484528y^4-270464y^3+137856y^2\\
&-19968y+256,\\
q_6(y)={}&16y^{10}-484y^9+5796y^8-34617y^7+105000y^6\\
&-135904y^5+7632y^4+60620y^3+57232y^2-52880y+64.
\end{align*}
Let \(V_f(x)\) denote the number of sign changes in the Sturm sequence of
\(f\) at \(x\). The endpoint variations are
\begin{center}
\begin{tabular}{c@{\qquad}ccc}
\toprule
factor & \(V_f(a_6)\) & \(V_f(b_6)\) & \(V_f(16)\)\\
\midrule
\(r_2\)     &0&0&0\\
\(r_{4,1}\)&2&2&2\\
\(r_{4,2}\)&0&0&0\\
\(p_6\)     &3&2&1\\
\(q_6\)     &2&2&1\\
\bottomrule
\end{tabular}
\end{center}
In particular, \(p_6\) has exactly one root in \(J_6\), while every other
factor in~\eqref{eq:app-resultant-factorization} has none there.  Since a
physical zero was established before the elimination, its square is the
unique root \(c_6\in J_6\) of \(p_6\).  This proves the algebraic
identification in Lemma~\ref{lem:c6-realization} without reversing the
logical implication from a resultant candidate to a physical state.

\subsection{A complete Grassmann atlas above \texorpdfstring{\(c_6\)}{c6}}
\label{appsubsec:g6-atlas}

We now cover
\[
 I=[b_6,16].
\]
The stable plane can be continued through the collision
\eqref{eq:app-repeat} without choosing the two multipliers separately.
Indeed, if \(t=P+P^{-1}\), then
\begin{equation}
 (t+2)(t-b(y))+a(y)^2=0.                              \label{eq:app-t-equation}
\end{equation}
The larger real solution \(t=t_+(y)>2\) selects the physical stable branch
by
\begin{equation}
 P=\frac{2}{t_+(y)+\sqrt{t_+(y)^2-4}},
 \qquad S=-\frac{a(y)P}{P+1}.                         \label{eq:app-SP-branch}
\end{equation}
The unstable plane uses the reciprocal multipliers.

For a cofactor row set \(abc\), let \(v_{abc}(z)\) be the signed vector of
the \(3\times3\) minors of the rows \(a,b,c\) of \(M-zI\).  Then
\[
 \Omega_{abc}(z_1,z_2)
 =\frac{v_{abc}(z_1)\wedge v_{abc}(z_2)}{z_1-z_2}
\]
is symmetric and extends at \(z_1=z_2=z\) to
\(v_{abc}(z)\wedge\partial_zv_{abc}(z)\).  This is the confluent wedge used
at \(y_*\).  A nonzero Pl\"ucker coordinate of \(\Omega_{abc}\) gives an
affine chart for the physical plane.

The following three closed charts cover \(I\):
\begin{center}
\small
\begin{tabular}{llll}
\toprule
interval & cofactor rows & right stable minor & left unstable minor\\
\midrule
\([b_6,31663/4000]\) & \(013\) & \(p_{23}\) & \(p_{01}\)\\
\([79157/10000,79159/10000]\) & \(012\) & \(p_{01}\) & \(p_{02}\)\\
\([158317/20000,16]\) & \(013\) & \(p_{23}\) & \(p_{01}\)\\
\bottomrule
\end{tabular}
\end{center}
For rows \(013\), elimination of either selected section against
\eqref{eq:app-product-relation} shows that its only possible zero in \(I\)
comes from
\[
 h(y)=3y^2-24y+2.
\]
Its Sturm variations at
\(7915780041490243/10^{15}\) and
\(7915780041490244/10^{15}\) are respectively one and zero.  Hence its
unique root lies strictly inside the middle chart and outside both outer
charts. On the middle interval, certified rational evaluation gives
\[
 p_{01}^{R}>0,\qquad p_{02}^{L}<0.
\]
The explicit margins are recorded in the supplement. Thus the bridge
sections do not vanish. The overlap intervals
\[
\left[\frac{79157}{10000},\frac{31663}{4000}\right],
\qquad
\left[\frac{158317}{20000},\frac{79159}{10000}\right]
\]
are nonempty, and the transition function
is the ratio of two sections already proved nonzero there.  Dividing by the
chosen section normalizes it to one.  This proves that the atlas has no
uncovered endpoint or transition point.

The repeated-multiplier energy \(y_*\) lies in the right outer chart.  The
confluent formula for \(\Omega_{013}\), together with the nonvanishing of
the selected section there, proves that the physical plane and its matching
determinant remain represented at \(y_*\).  Equivalently, the root of
\(12y^2-96y-17\) in \(I\) is a regular point of the Grassmann curve even
though the individual multiplier labels meet.

\subsection{Exhaustion and genuine unsquared exclusion}
\label{appsubsec:g6-exhaustion}

Every zero of the physical determinant on \(I\) must annihilate the
resultant~\eqref{eq:app-resultant-factorization}.  The Sturm table above
shows that only \(p_6\) and \(q_6\) have roots in \(I\), one each.  They are
isolated by
\begin{align}
 I_1={}&\left[
 \frac{8080985802104273}{10^{15}},
 \frac{4040492901052137}{5\cdot10^{14}}
 \right],                                             \label{eq:app-I1}\\
 I_2={}&\left[
 \frac{406992828166963}{5\cdot10^{13}},
 \frac{101748207041741}{12500000000000}
 \right].                                             \label{eq:app-I2}
\end{align}
There are no other resultant roots on \(I\).  The chart-section root near
\(7.915780041490243\) is not in this list because the bridge chart shows it
to be a coordinate degeneracy only.  The repeated-multiplier value
\eqref{eq:app-repeat} is represented by the confluent quotient and is not a
missing resultant root.

For cross-chart necessity, at every
\(y\in I\), the atlas provides a chart in which both selected Pl\"ucker
sections are nonzero.  In that chart, the normalized numerator differs from
the coordinate-free exterior determinant only by the product of those
nonzero sections and the nonzero ordered Vandermonde factors.  Hence a
physical zero makes the chart numerator vanish.  Clearing only the
denominators already proved nonzero gives
\[
R(P,y)=G(P,y)=0,
\]
so the resultant must vanish.  At the sole section zero the middle chart
supplies this implication; at \(y_*\) the confluent wedge replaces the
Vandermonde quotient and gives the same implication. Consequently the
candidate list covers chart transitions and the repeated multiplier.

The final exclusion uses the original unsquared equation. In chart \(013\),
substitution of the physical branch
\eqref{eq:app-SP-branch} into~\eqref{eq:app-Evans}, followed by rational
interval evaluation, gives
\[
\sup_{y\in I_1}\mathcal E_6(y)<0,
\qquad
\sup_{y\in I_2}\mathcal E_6(y)<0.
\]
The corresponding outward enclosures are recorded in the supplement.
Reconstructing the same planes from
cofactor rows \(023\) gives strictly negative rational enclosures on
\(I_1\) and \(I_2\) as well. The latter reconstruction uses different
Pl\"ucker sections and provides a second normalization check.

The root in \(I_1\) is a physical root for a different defect transfer, but
the first strict sign bound shows that it is not a G6 match. The root in
\(I_2\) comes from the second algebraic branch retained by squaring and
clearing denominators; the second strict sign bound excludes it as well.
Consequently the positive G6 determinant has no zero on \(I\).  Together
with the unique physical root in \(J_6\), the essential-spectrum reduction,
and \(\|A_6\|\leq4\), this proves that \(+\sqrt{c_6}\) is the largest
positive spectral point of \(A_6\).

\subsection{The symmetry and the squared multiplicity}
\label{appsubsec:g6-rank-two}

We finish by auditing the negative branch and the multiplicity after
squaring.  Let
\[
 D_6^{+}=\{-4j:j\geq0\}\cup\{6+4j:j\geq0\}
\]
be the set of positive \(Q\)-sites.  The reflection \(i\mapsto6-i\)
exchanges its two displayed parts, so
\begin{equation}
 Q_{6-i}=Q_i\qquad(i\in\mathbb Z).                    \label{eq:app-Q-reflect}
\end{equation}
The finite anchor computation
\[
 (Q_0,\ldots,Q_6)=(1,-1,-1,-1,-1,-1,1),\qquad
 (\tau_0,\ldots,\tau_7)=(1,1,-1,1,-1,1,-1,-1)
\]
gives \(\tau_7=-\tau_0\).  If \(\tau_{7-i}=-\tau_i\), then
\[
 \tau_{6-i}=Q_{6-i}\tau_{7-i}
 =-Q_i\tau_i=-\tau_{i+1}.
\]
The same recurrence run backwards proves
\begin{equation}
 \tau_{7-i}=-\tau_i\qquad(i\in\mathbb Z).             \label{eq:app-tau-reflect}
\end{equation}

Define the unitary operator
\[
 (Ku)_i=(-1)^iu_{9-i}.
\]
Then
\[
 (K^2u)_i=(-1)^i(-1)^{9-i}u_i=-u_i,
\]
so \(K^2=-I\).  Moreover,
\begin{align*}
(KA_6u)_i=(-1)^i\bigl(&u_{8-i}+u_{10-i}
 +\tau_{7-i}u_{7-i}+\tau_{9-i}u_{11-i}\bigr),\\
(A_6Ku)_i=(-1)^i\bigl(&-u_{10-i}-u_{8-i}
 +\tau_{i-2}u_{11-i}+\tau_i u_{7-i}\bigr).
\end{align*}
Using~\eqref{eq:app-tau-reflect} at \(i\) and \(i-2\) gives
\[
 KA_6=-A_6K,\qquad KH_6=H_6K.                         \label{eq:app-K-relations}
\]
Thus \(K\) maps the one-dimensional eigenspace at \(+\sqrt{c_6}\)
bijectively onto a one-dimensional eigenspace at \(-\sqrt{c_6}\).
It also maps any hypothetical negative spectral point of larger absolute
value to a positive point already excluded above.

Finally, because \(A_6\) is self-adjoint and \(c_6>0\),
\begin{equation}
 \ker(H_6-c_6)
 =\ker(A_6-\sqrt{c_6})
 \mathbin{\overset{\perp}{\oplus}}
 \ker(A_6+\sqrt{c_6}).                                \label{eq:app-rank-two}
\end{equation}
Both summands have dimension one.  Therefore
\[
 \sup\sigma(H_6)=c_6,\qquad
 \dim\ker(H_6-c_6)=2.
\]
Reflection reverses the interface by unitary conjugacy, and negating the
lift gives a diagonally conjugate squared operator.  Hence the same global
edge and rank-two conclusion hold for both orientations and both lifts.
